\documentclass[11pt]{article}
\newcommand{\documenttype}{Document}
\usepackage[a4paper,margin=0.85in]{geometry}
\usepackage[T1]{fontenc}
\usepackage[utf8]{inputenc}
\usepackage{lmodern,amsmath,amssymb,microtype,xcolor,enumitem,tabularx,titlesec,fancyhdr,hyperref}
\definecolor{accent}{HTML}{315B8A}
\definecolor{link}{HTML}{5B4F7A}
\hypersetup{colorlinks=true,urlcolor=link,linkcolor=link,citecolor=link,pdfauthor={Phurailatpam Hemantakumar}}
\setlength{\parindent}{0pt}\setlength{\parskip}{5pt}\setlength{\headheight}{14pt}\setlength{\emergencystretch}{2em}
\titleformat{\section}{\large\bfseries\color{accent}}{}{0pt}{}
\titlespacing*{\section}{0pt}{10pt}{5pt}
\setlist[itemize]{leftmargin=16pt,itemsep=2pt,topsep=3pt,parsep=0pt}
\pagestyle{fancy}\fancyhf{}\fancyhead[L]{\small Phurailatpam Hemantakumar}\fancyhead[R]{\small\documenttype}\fancyfoot[R]{\small\thepage}\renewcommand{\headrulewidth}{0.3pt}
\newcommand{\firstheading}[2]{\thispagestyle{plain}{\LARGE\bfseries\color{accent}#1}\par\vspace{3pt}{#2}\par\vspace{4pt}{\color{accent}\hrule}\vspace{8pt}}
\newcommand{\dated}[2]{\par\noindent\begin{tabularx}{\linewidth}{@{}X r@{}}#1&#2\end{tabularx}\par}
\newcommand{\contact}{\href{mailto:hemantaph@link.cuhk.edu.hk}{hemantaph@link.cuhk.edu.hk}\quad|\quad +852 9867 8178\quad|\quad\href{https://hemantaph.com}{hemantaph.com}\\Department of Physics, The Chinese University of Hong Kong}
\renewcommand{\documenttype}{Cover letter}
\hypersetup{pdftitle={Cover letter | Monash OzGrav Research Fellow 698908}}
\setlength{\parskip}{7pt}
\begin{document}
\firstheading{Phurailatpam Hemantakumar}{\contact}
21 November 2026

Selection Committee\\
Research Fellow -- Gravitational Wave Discovery (OzGrav), Job No. 698908\\
School of Physics and Astronomy, Monash University\\
Clayton, Victoria, Australia

Dear Members of the Selection Committee,

I am applying for the Research Fellow position in Gravitational Wave Discovery (OzGrav), Job No. 698908, at Monash University. I have completed my PhD research and examination in Physics at The Chinese University of Hong Kong under Otto Akseli Hannuksela, with the formal award expected in December 2026. My research focuses on strong gravitational-wave (GW) lensing statistics, compact-binary populations, Bayesian data analysis, and eccentric waveform modelling. I am a member of the LIGO Scientific Collaboration and have developed scientific software used for GW event-rate and detectability calculations.

I lead development of \href{https://ler.hemantaph.com/}{\texttt{ler}}, an open-source Monte Carlo framework for simulating lensed and unlensed compact-binary populations and their detectable rates (\textit{JOSS}, 2026). My work spans population modelling, elliptical power-law lens models with external shear, numerical cross-sections, optical depths, detector selection, and the propagation of model uncertainties. I also lead development of \href{https://gwsnr.hemantaph.com/}{\texttt{gwsnr}}, which provides fast signal-to-noise ratio and detection-probability calculations using parallelised waveform inner products and accelerated statistical methods. These tools have supported GW lensing follow-up analyses and provide a basis for population inference as detector sensitivity improves.

In complementary research, I co-developed \texttt{TaylorF2Ecck}, a frequency-domain post-Newtonian eccentric inspiral approximant. I implemented the model in a \texttt{LALSuite} fork and carried out Bayesian parameter estimation for GW170817 and GW190425 (\textit{Physical Review D}, 2026). I have also contributed to studies of eccentricity and line-of-sight acceleration, and to multimessenger forecasts for lensed GW and short gamma-ray burst counterparts.

The breadth of GW research at Monash's OzGrav node is particularly attractive. I would contribute to LVK data analysis and compact-binary population studies through consistent modelling of merger-rate histories, detector selection, and strong-lensing probabilities. My planned work on star-formation-informed binary-black-hole merger rates and on testing population-based lensing predictions is complementary to the node's strengths in GW astrophysics and Bayesian inference. I would also welcome opportunities to connect eccentric waveform systematics with broader studies of compact binaries and tests of gravity.

I bring experience leading collaborative software projects, publishing refereed research, mentoring students, and teaching computational methods. I would value the opportunity to develop these research directions with Professors Eric Thrane, Paul Lasky, Ilya Mandel, and the wider OzGrav community. I enclose my CV, response to the key selection criteria, and supporting research materials.

Thank you for considering my application.

Sincerely,\\
Phurailatpam Hemantakumar
\end{document}
